\section{微分概念及其计算}\label{sec:4.2}

微商的概念，描述了函数在一点的变化率，它自然地也就描述了函数在一点附近的变化情况．考虑函数 $y=f(x)$ 当自变量 $x$ 的改变量为 $\Delta x$ 时，函数相应的改变量 $\Delta y$ 的变化．

设 $y=f(x)$ 在 $x$ 点可导，即下面的极限存在：
\[
f'(x)=\lim_{\Delta x\to 0}\frac{\Delta y}{\Delta x}=\lim_{\Delta x\to 0}\frac{f(x+\Delta x)-f(x)}{\Delta x}.
\]
因此
\[
\frac{\Delta y}{\Delta x}=f'(x)+\alpha,
\]
其中 $\alpha$ 是一个无穷小量（当 $\Delta x\to 0$），即 $\alpha\to0$ ($\Delta x\to0$)．于是
\[
\Delta y=f'(x)\Delta x+\alpha\Delta x=f'(x)\Delta x+o(\Delta x).
\]
这样 $\Delta y$ 被分解为两项之和，第一项和 $\Delta x$ 成正比，也就是 $\Delta x$ 的线性函数；第二项是当 $\Delta x\to0$ 时比 $\Delta x$ 更高阶的无穷小量．如果把 $y=f(x)$ 视为时间 $x$ 时所走过的路程，那么 $f'(x)$ 就是时刻 $x$ 时的运动速度．于是第一项 $f'(x)\Delta x$ 是以匀速 $f'(x)$ 运动，在 $\Delta x$ 时间内所走过的路程；而第二项则是因为加速度的作用而产生的附加路程．尽管这里没有写出第二项的具体表达式，但我们知道它是比 $\Delta x$ 高阶的无穷小量．也就是说，当 $|\Delta x|$ 很小时，它比第一项要小得多，因此记第一项为 $\Delta y$ 的主要部分．根据这样一些特性，我们给出如下定义．

\begin{definition}\label{def:4-2}
设 $y=f(x)$ 在 $(a,b)$ 有定义，如果对给定的 $x\in(a,b)$，有
\[
\Delta y=f(x+\Delta x)-f(x)=A\Delta x+o(\Delta x)\quad(\Delta x\to0),
\]
其中 $A$ 与 $\Delta x$ 无关，则称 $f(x)$ 在 $x$ 点可微，并称 $A\Delta x$ 为函数 $f(x)$ 在 $x$ 点的微分，记为
\[
\mathrm{d}y=A\Delta x \quad\text{或}\quad \mathrm{d}f(x)=A\Delta x.
\]
\end{definition}

上述定义中有两个概念，一个是可微的概念，另一个是微分的概念．从定义可知，微分具有两大重要特性：
\begin{enumerate}[1)]
\item 微分是自变量的改变量 $\Delta x$ 的线性函数；
\item 微分与函数的改变量 $\Delta y$ 之差是比 $\Delta x$ 高阶的无穷小量．
\end{enumerate}

由于特性 2)，我们称微分 $\mathrm{d}y$ 为改变量 $\Delta y$ 的主要部分，又由于特性 1)，微分 $\mathrm{d}y$ 进一步称为改变量 $\Delta y$ 的线性主要部分．事实上当 $\mathrm{d}y\neq0$ 时
\[
\begin{aligned}
\lim_{\Delta x\to0}\frac{\Delta y}{\mathrm{d}y}&=\lim_{\Delta x\to0}\frac{\mathrm{d}y+o(\Delta x)}{\mathrm{d}y}\\
&=\lim_{\Delta x\to0}\left(1+\frac{o(\Delta x)}{A\Delta x}\right)=1\text{，}
\end{aligned}
\]
即 $\Delta y$ 与 $\mathrm{d}y$ 是等价无穷小量．

显然，当 $x$ 在 $(a,b)$ 中变化时，系数 $A$ 也随 $x$ 变化，因此它是 $x$ 的函数，故微分 $\mathrm{d}y$ 既与 $x$ 有关，又与 $\Delta x$ 有关，而 $x$ 和 $\Delta x$ 是两个互相独立的变量，但它对 $\Delta x$ 的依赖是线性的．

\begin{example}\label{ex:4-2-1}
自由落体运动中，物体下落的路程 $s$ 与时间 $t$ 的关系是
\[
s(t)=\frac{1}{2}gt^{2}.
\]
根据
\[
\begin{aligned}
\Delta s&=s(t+\Delta t)-s(t)=\frac{1}{2}g(t+\Delta t)^{2}-\frac{1}{2}gt^{2}\\
&=\frac{1}{2}g\left(2t\Delta t+(\Delta t)^{2}\right)=gt\Delta t+\frac{1}{2}g(\Delta t)^{2},
\end{aligned}
\]
即 $\Delta s$ 可表为 $\Delta t$ 的线性函数和 $\Delta t$ 的高阶无穷小量之和，由微分定义知，$s(t)$ 在 $t$ 点可微，且微分
\[
\mathrm{d}s=gt\Delta t.
\]
它等于我们一开始讲过的，以匀速 $s'(t)=gt$ 运动，在 $\Delta t$ 时间内所走过的路程．
\end{example}

\begin{example}\label{ex:4-2-2}
圆面积 $y=\pi R^{2}$，在 $R$ 有改变量 $\Delta R$，对应地函数有改变量
\[
\begin{aligned}
\Delta y&=\pi(R+\Delta R)^{2}-\pi R^{2}\\
&=2\pi r\Delta R+\pi(\Delta R)^{2}.
\end{aligned}
\]
$\Delta y$ 可表示为 $\Delta R$ 的线性函数与 $\Delta R$ 的高阶无穷小之和，故函数在 $R$ 可微，且微分
\[
\mathrm{d}y=2\pi R\Delta R.
\]
\end{example}

从几何上（图~\ref{fig:4-5}）看，$\pi R^{2}$ 表示半径为 $R$ 的圆面积．微分可以这样理解，当半径 $R$ 变大即圆面积膨胀时，设想圆周长保持不变，半径增大 $\Delta R$ 所引起的圆面积变化 $2\pi R\cdot\Delta R$，这就是圆面积的微分．它与 $\Delta R$ 成正比，与圆面积真正的变化之差是较 $\Delta R$ 高阶的无穷小．当然圆不可能保持周长不变而膨胀，这只是一种设想而已．但当 $\Delta R$ 很小时，两者之差就更小了．

\begin{example}\label{ex:4-2-3}
设正方形的边长为 $x$，则面积为
\[
f(x)=x^{2}.
\]
\[
\begin{aligned}
\Delta f(x)&=(x+\Delta x)^{2}-x^{2}\\
&=2x\Delta x+(\Delta x)^{2},
\end{aligned}
\]
即 $\Delta f(x)$ 可表为 $\Delta x$ 的线性函数和 $\Delta x$ 的高阶无穷小量之和，故 $f(x)$ 在 $x$ 点可微，且微分
\[
\mathrm{d}y=2x\Delta x.
\]
\end{example}

从几何上看（图~\ref{fig:4-6}），$2x\Delta x$ 即是图中阴影部分的面积，而 $(\Delta x)^{2}$ 便是以 $\Delta x$ 为边长的小正方形的面积．当边长 $x$ 有一个微小增量 $\Delta x$ 时，面积的增量可近似地用阴影部分代替，而忽略的部分仅是长为 $\Delta x$ 的小正方形面积．和圆面积的情形类似，同样可以设想边长为 $x$ 的正方形，当 $x$ 变大，即正方形膨胀时，正方形的总边长保持不变所产生的面积变化 $2x\Delta x$，即为面积的微分．

\begin{figure}[htbp]
\centering
\begin{minipage}[t]{0.45\textwidth}
\centering
\includegraphics[width=0.8\textwidth]{fig4-5.png}
\figcap
\label{fig:4-5}
\end{minipage}
\begin{minipage}[t]{0.45\textwidth}
\centering
\includegraphics[width=0.8\textwidth]{fig4-6.png}
\figcap
\label{fig:4-6}
\end{minipage}
\end{figure}

上面三个例子中的函数都很简单，可以想象，若函数稍微复杂一点，按可微的定义去求 $\Delta y$ 的分解式，并不是件容易的事．其实，在本节的开头，我们已推出：若 $f(x)$ 在 $x$ 可导，则 $f(x)$ 在 $x$ 可微，且 $\mathrm{d}y=f'(x)\Delta x$．进一步我们有下面的定理．

\begin{theorem}\label{thm:4-5}
函数 $y=f(x)$ 在 $x$ 点可微的充要条件是：函数 $f(x)$ 在 $x$ 点可导．这时微分中 $\Delta x$ 的系数 $A=f'(x)$．
\end{theorem}

\begin{proof}
充分性前面已证，下证必要性．设 $f(x)$ 在 $x$ 点可微，由定义知
\[
\Delta y=A\Delta x+o(\Delta x),
\]
因此
\[
\lim_{\Delta x\to0}\frac{\Delta y}{\Delta x}=\lim_{\Delta x\to0}\frac{A\Delta x+o(\Delta x)}{\Delta x}=A.
\]
故 $f(x)$ 在 $x$ 点可导，且 $f'(x)=A$．
\end{proof}

这个定理告诉我们，对一元函数而言，可微与可导是等价的，且有关系式
\[
\mathrm{d}y=f'(x)\Delta x.
\]
这个公式使我们可以方便地求出函数的微分．如前面三例中：
\[
s(t)=\frac{1}{2}gt^{2},\ s'(t)=gt,\quad\text{则}\ \mathrm{d}s(t)=gt\Delta t.
\]
\[
f(R)=\pi R^{2},\ f'(R)=2\pi R,\ \text{则}\ \mathrm{d}f(R)=2\pi R\Delta R.
\]
\[
f(x)=x^{2},\ f'(x)=2x,\ \text{则}\ \mathrm{d}f(x)=2x\Delta x.
\]
若 $y=x$，则 $y'=1$，于是
\[
\mathrm{d}y=\mathrm{d}x=1\cdot\Delta x=\Delta x.
\]
因此我们规定自变量的微分 $\mathrm{d}x$ 等于自变量的改变量 $\Delta x$．这样微分公式又可写为 $\mathrm{d}y=f'(x)\mathrm{d}x$，于是有 $\dfrac{\mathrm{d}y}{\mathrm{d}x}=f'(x)$．在定义微商时，符号 $\dfrac{\mathrm{d}y}{\mathrm{d}x}$ 是作为一个整体，而现在有了微分概念之后，微商可以看作是微分之商．也就是说，微商的确是微分之商．

历史上，牛顿把自变量理解为时间 $t$，称函数 $x$ 为流量，称函数的变化速度（变化率）为流数，记作 $\dot{x}$，把函数的无穷小改变称为瞬（moment），记为 $x0$，这就是他的直观的微分概念．对他来说，微商是微分之商是理所当然的．由于没有清楚的极限概念，他的微商微分概念，都是不精确的．后来有了极限概念后，才把微商概念搞清楚，并给出了微分是函数改变量的线性主要部分的定义，从而也得到了微商是微分之商这个理论与直观统一的结论：$\mathrm{d}y=f'(x)\mathrm{d}x$．根据这一点，历史上微商也有称为微分系数的（differential coefficient）．

借助于微商的几何意义，我们立即可得出

微分的几何意义．由于
\[
\mathrm{d}y=f'(x)\mathrm{d}x=\tan\alpha\,\mathrm{d}x,
\]
因此微分 $\mathrm{d}y$ 是曲线 $y=f(x)$ 在 $(x,y)$ 处的切线对应的改变量．用微分 $\mathrm{d}y$ 近似地代替改变量 $\Delta y$，从几何上看就是用切线的改变量近似地代替函数的改变量（图~\ref{fig:4-7}）．

\begin{figure}[htbp]
\centering
\includegraphics[width=0.4\textwidth]{fig4-7.png}
\figcap
\label{fig:4-7}
\end{figure}

从微商表，我们很容易得到基本初等函数的微分公式，如
\[
\mathrm{d}x^{\alpha}=\alpha x^{\alpha-1}\mathrm{d}x\text{；}\qquad
\mathrm{d}e^{x}=e^{x}\mathrm{d}x\text{；}\qquad
\mathrm{d}\ln x=\frac{1}{x}\mathrm{d}x\text{；}\qquad
\mathrm{d}\sin x=\cos x\mathrm{d}x
\]
等等．同样借助于微商的运算法则，立即可得下面的微分运算法则：

（1）四则运算法则．
\[
\mathrm{d}(f(x)\pm g(x))=\mathrm{d}f(x)\pm\mathrm{d}g(x)\text{；}
\]
\[
\mathrm{d}(f(x)g(x))=g(x)\mathrm{d}f(x)+f(x)\mathrm{d}g(x)\text{；}
\]
\[
\mathrm{d}\left(\frac{f(x)}{g(x)}\right)=\frac{g(x)\mathrm{d}f(x)-f(x)\mathrm{d}g(x)}{g^{2}(x)}\quad(g(x)\neq0)\text{．}
\]

（2）复合函数的微分．设 $y=f(u)$，$u=g(x)$，则复合函数 $y=f(g(x))=\varphi(x)$ 的微分为
\[
\mathrm{d}y=(f\circ g)'(x)\mathrm{d}x=f'(g(x))g'(x)\mathrm{d}x.
\]
上式中，实际上 $f'(g(x))=f'(u)$，$g'(x)\mathrm{d}x=\mathrm{d}u$，因此微分 $\mathrm{d}y$ 也可写成 $\mathrm{d}y=f'(u)\mathrm{d}u$．

把 $\mathrm{d}y=f'(u)\mathrm{d}u$ 与 $\mathrm{d}y=\varphi'(x)\mathrm{d}x$ 相比较，虽然 $x$ 是自变量，$u$ 是中间变量，但两者形式上是一样的，这一性质称为\emph{一阶微分形式的不变性}．一阶微分形式不变性说明，可以在微分等式中代入变量．例如 $y=e^{u}$，$u=x^{2}$，则
\[
\mathrm{d}y=e^{u}\mathrm{d}u,
\]
代入变量 $u=x^{2}$ 得
\[
\mathrm{d}y=e^{x^{2}}\mathrm{d}x^{2}=e^{x^{2}}\cdot 2x\mathrm{d}x.
\]
这种``代入''运算，在微商公式中就不可以做．例如在 $y'=e^{u}$ 中代入变量 $u=x^{2}$，得 $y'=e^{x^{2}}$，显然结果是错误的．以后我们将看到，正是由于在微分式中可以任意作变量代入，使得许多时候微分运算比微商运算更重要．

\begin{example}\label{ex:4-2-4}
设
\[
y=e^{\sin(ax+b)}+\frac{x}{1+x^{2}}.
\]
利用微分运算法则求函数的微分．

解
\[
\begin{aligned}
\mathrm{d}y&=e^{\sin(ax+b)}\mathrm{d}\sin(ax+b)+\frac{(1+x^{2})\mathrm{d}x-x\mathrm{d}(1+x^{2})}{(1+x^{2})^{2}}\\
&=e^{\sin(ax+b)}\cos(ax+b)\mathrm{d}(ax+b)+\frac{(1+x^{2})\mathrm{d}x-2x^{2}\mathrm{d}x}{(1+x^{2})^{2}}\\
&=e^{\sin(ax+b)}a\cos(ax+b)\mathrm{d}x+\frac{1-x^{2}}{(1+x^{2})^{2}}\mathrm{d}x\\
&=\left[e^{\sin(ax+b)}a\cos(ax+b)+\frac{1-x^{2}}{(1+x^{2})^{2}}\right]\mathrm{d}x.
\end{aligned}
\]
\end{example}

\begin{example}\label{ex:4-2-5}
求 $\sin 31^{\circ}$ 的近似值．

解 令 $f(x)=\sin x$，$x=30^{\circ}=\dfrac{\pi}{6}$，$\Delta x=1^{\circ}=\dfrac{\pi}{180}$．用近似公式 $\Delta f(x)\approx\mathrm{d}f(x)$ 得
\[
f(x+\Delta x)\approx f(x)+f'(x)\Delta x,
\]
即
\[
\begin{aligned}
\sin 31^{\circ}&=\sin\left(\frac{\pi}{6}+\frac{\pi}{180}\right)\approx\sin\frac{\pi}{6}+\cos\frac{\pi}{6}\cdot\frac{\pi}{180}\\
&\approx\frac{1}{2}+\frac{\sqrt{3}}{2}\times0.017\,45=0.515\,1.
\end{aligned}
\]
\end{example}

\begin{xiti}
\item 求下列函数在指定点的微分：
\begin{ziti}
\item $y=a_n x^n + a_{n-1}x^{n-1} + \cdots + a_1 x + a_0$，求 $\mathrm{d}y(0)$，$\mathrm{d}y(1)$；
\item $y = \sec x + \tan x$，求 $\mathrm{d}y(0)$，$\mathrm{d}y\left(\dfrac{\pi}{4}\right)$ 和 $\mathrm{d}y(\pi)$；
\item $y = \dfrac{1}{a}\arctan\dfrac{x}{a}$，求 $\mathrm{d}y(0)$，$\mathrm{d}y(a)$；
\item $y = \dfrac{1}{x} + \dfrac{1}{x^2}$，求 $\mathrm{d}y(0.1)$，$\mathrm{d}y(0.01)$．
\end{ziti}

\item 求下列函数的微分：
\begin{ziti}
\item $y = \dfrac{x}{1-x^2}$；
\item $y = x\ln x - x$；
\item $y = \sqrt{x} + \ln x - \dfrac{1}{\sqrt{x}}$；
\item $y = \arcsin\sqrt{1-x^2}$；
\item $y = e^{\sin x^2}$；
\item $y = \ln\left|\tan\left(\dfrac{x}{2} + \dfrac{\pi}{4}\right)\right|$．
\end{ziti}

\item 设 $u$，$v$ 是 $x$ 的可微函数，求 $\mathrm{d}y$：
\begin{ziti}
\item $y = \arctan\dfrac{u}{v}$；
\item $y = \ln\sqrt{u^2 + v^2}$；
\item $y = \ln\sin(u + v)$；
\item $y = \dfrac{1}{\sqrt{u^2 + v^2}}$．
\end{ziti}

\item 求下列函数的微分 $\mathrm{d}y$：
\begin{ziti}
\item $y = \sin^2 t$，$t = \ln(3x + 1)$；
\item $y = \ln(3t + 1)$，$t = \sin^2 x$；
\item $y = e^{3u}$，$u = \dfrac{1}{2}\ln t$，$t = x^3 - 2x + 5$；
\item $y = \arctan u$，$u = (\ln t)^2$，$t = 1 + x^2 - \cot x$．
\end{ziti}

\item 求下列各式的近似值：
\begin{ziti}
\item $\sqrt{120}$；
\item $\arctan 1.05$；
\item $\sin 29^\circ$；
\item $\sqrt{\dfrac{(2.037)^2 - 1}{(2.037)^2 + 1}}$．
\end{ziti}
\end{xiti}
